\section{本章小结}
\label{sec:metric-summary}

本章从实数的绝对值、向量的范数、二进制串的Hamming距离和图的最短路距离中提炼出度量的三条要求，并说明度量是空间结构的一部分：同一集合换一个度量，球的形状、哪些点列收敛、是否可分都可能改变。距离依次提供了三层语言。开球给出邻近结构，由它定义开集、闭集、闭包与稠密；点列收敛把邻近变成可以逐项检查的误差，闭集的序列判据把两者连起来；连续映射研究误差在两个空间之间的传递，一致连续与Lipschitz条件分别要求统一的容差与统一的放大倍数。

几个“有限性”与“存在性”条件分别控制不同的对象。有界控制集合的总体范围；可分要求一套可数代表服务所有精度；全有界要求每个精度只需有限个代表；完备要求柯西列在空间内有极限；紧致要求每个给定的开覆盖都有有限子覆盖。在度量空间中，紧致等价于序列紧致，也等价于完备且全有界；只有在$\mathbb R^d$这样的有限维欧氏空间中，才可以简化为闭且有界。紧致定义域使连续函数取得最值、使最近点存在、使连续升级为一致连续，但不提供Lipschitz常数。比较两个空间时，双Lipschitz映射保留柯西列、完备性和全有界性，同胚只保留开集与收敛；logit映射表明，完备性不是由开集单独决定的。

压缩映射定理把这些结论串成一条完整的论证：统一的压缩常数使相邻差几何衰减，几何尾和使迭代列成为柯西列，完备性给出极限，连续性验证极限满足$T(x^\star)=x^\star$，压缩条件再给出唯一性。同一条尾和估计给出先验与后验误差界，带误差的版本说明执行误差会留下$\eta/(1-q)$量级的偏差。各个反例提醒我们：只观察到相邻项越来越近，或在局部看到误差缩小，都不足以断定极限存在或结果准确；每个结论都需要它的全部条件。

下一章转向另一个问题：怎样测量一个集合的大小。度量空间的开集将生成Borel集合族，为测量提供可用的集合；但距离本身并不决定集合的测度，同一个度量空间上可以有许多不同的测度。
